%====================== CONSIDERAÇÕES FINAIS ======================
\section{Considerações finais}

Este trabalho apresentou o desenvolvimento do LibrasAI, um sistema de
reconhecimento automático de gestos dinâmicos de Libras em tempo real, construído
a partir de um \textit{pipeline} que integra captura de vídeo, extração de
\textit{landmarks} via MediaPipe, normalização espacial e classificação por meio
de uma rede híbrida CNN + LSTM.

Os objetivos específicos definidos para este estágio do trabalho foram
majoritariamente alcançados: o \textit{pipeline} de captura e detecção de
\textit{landmarks} está implementado e em funcionamento; o modelo CNN + LSTM foi
treinado e superou a meta de acurácia estabelecida (90\%), atingindo 97,4\% sobre
o \textit{dataset} inicial; o sistema de coleta interativa de dados está
operacional e já produziu um protocolo replicável com 10 participantes; e a
formação automática de frases, com os três modos de operação planejados, foi
implementada e testada pela equipe de desenvolvimento.

Ao mesmo tempo, este trabalho reconhece explicitamente que os resultados obtidos
são preliminares e devem ser lidos à luz das limitações discutidas na
Seção~\ref{subsec:limitacoes} — em especial o tamanho ainda reduzido do
\textit{dataset} e a necessidade de validação com usuários externos ao
treinamento, endereçada pelo protocolo descrito na Seção~\ref{subsec:validacao}.
Essas limitações não invalidam a contribuição do trabalho, mas delimitam com
precisão o que ainda precisa ser feito antes que o sistema possa ser considerado
robusto o suficiente para uso real pela comunidade surda.

Do ponto de vista de contribuição acadêmica, o LibrasAI se soma a um conjunto
ainda pequeno de sistemas de reconhecimento de Libras — em contraste com a
abundância de trabalhos voltados à ASL e a outras línguas de sinais
internacionais —, reforçando a relevância de pesquisa nacional dedicada
especificamente à língua de sinais brasileira. A intenção declarada de dar
continuidade a este trabalho em nível de mestrado, com submissão dos resultados a
um artigo científico, indica que o presente TCC deve ser compreendido como uma
etapa inicial de um programa de pesquisa mais amplo, e não como um produto
finalizado.
